\section{Graph Energy Matching}
\label{sec:gem}
We extend continuous Energy Matching methods to discrete graph spaces through a novel transport-aligned discrete proposal tailored specifically for discrete probability distributions. Our formulation employs an energy-based switching mechanism that clearly separates \emph{transport} and \emph{mixing} sampling regimes without requiring explicit indexing of time or noise levels. Sampling is efficiently facilitated through deterministic greedy proposals in the transport regime or stochastic gradient-informed proposals during mixing, alongside an annealing strategy to efficiently generate novel samples from optionally given data samples via incremental edits. The sampler uses a temperature $\epsilon$ to control stochasticity: the transport regime is the deterministic $\epsilon\to 0$ limit, while the mixing regime uses $\epsilon>0$.

\paragraph{Graph representation.} 
We consider a discrete state space $\mathcal{X}:=\big(\mathcal{X}_{\mathrm{node}}\big)^{n}\times\big(\mathcal{X}_{\mathrm{edge}}\big)^{\binom{n}{2}}$ to represent undirected graphs. A graph $x\in \mathcal{X}$ contains a set of $n$ nodes $\mathcal{V}=\{1,\dots,n\}$ and edges $\mathcal{E}$ defining the connectivity between these nodes. Each node $i\in\mathcal{V}$ has a node class $x_i \in \mathcal{X}_{\mathrm{node}}$ (e.g.\ an atom type), and each edge $(i,j)\in\mathcal{E}$ has an edge class $x_{ij}\in\mathcal{X}_{\mathrm{edge}}$ (e.g.\ a bond type). The edge space $\mathcal{X}_{\mathrm{edge}}$ includes an element indicating the absence of an edge.




The graph $x:=\big(x_\mathcal{V},x_\mathcal{E}\big)$ consists of $x_\mathcal{V}$ and $x_\mathcal{E}$ denoting categorical node and edge features, respectively, and are defined as
\begin{equation}
x_\mathcal{V}  :=[x_i]_{i=1}^n\in(\mathcal{X}_{\mathrm{node}})^n;\quad
x_\mathcal{E}  :=[x_{ij}]_{ i<j}\in\big(\mathcal{X}_{\mathrm{edge}}\big)^{\binom{n}{2}}\nonumber
\end{equation}

We embed the node and edge categorical feature vectors $x_\mathcal{V}$ and $x_\mathcal{E}$ using a real-valued one-hot encoding and treat them as continuous features via
\[
\phi_{\mathcal{V}} :\mathcal{X}_{\mathrm{node}}\to\mathbb{R}^{l_{\mathrm{node}}};\quad
\phi_{\mathcal{E}} :\mathcal{X}_{\mathrm{edge}}\to\mathbb{R}^{l_{\mathrm{edge}}}
\]
where $l_{\mathrm{node}}=|\mathcal{X}_{\mathrm{node}}|$ and $l_{\mathrm{edge}}=|\mathcal{X}_{\mathrm{edge}}|$ denote the numbers of node and edge classes, respectively. We denote the real-valued node features $\hat{x}_{\mathcal{V}}\in\mathbb{R}^{n l_{\mathrm{node}}}$ and edge features $\hat{x}_{\mathcal{E}}\in\mathbb{R}^{\binom{n}{2}l_{\mathrm{edge}}}$, and the corresponding graph representation 
$\hat{x}\in\mathbb{R}^d$, where $d=n l_{\mathrm{node}}+\binom{n}{2}l_{\mathrm{edge}}$, as:
\begin{equation}
\hat{x}_\mathcal{V}:=\oplus_{i=1}^n \phi_\mathcal{V}(x_i);\quad
\hat{x}_\mathcal{E}:=\oplus_{i<j}\phi_\mathcal{E}(x_{ij});\quad
\hat{x}:=\hat{x}_\mathcal{V}\oplus \hat{x}_\mathcal{E}. \nonumber
\end{equation}
where $\oplus$ denotes the element-wise concatenation operation.
We parameterize the energy as $V_\theta(x):=\tilde V_\theta(\hat{x})$, differentiable with respect to
the continuous embedding $\hat{x}$.

\paragraph{Local and Permutation-Invariant Geometries.} The potential $V_\theta$ is permutation-invariant by construction (Appendix~\ref{app:perm-invariance}), so the geometry used for pairing should also respect node relabelings. Here, $\sigma\in S_n$ acts on node indices and induces a relabeling of the graph: $(\sigma\cdot x)_i=x_{\sigma(i)}$ and $(\sigma\cdot x)_{ij}=x_{\sigma(i)\sigma(j)}$. We measure local discrepancies in the embedding by
\begin{equation}\label{eq:loc-cost}
c_{\mathrm{loc}}(x,y)
\;\coloneqq\;
\lambda_\mathcal{V}\|\hat{x}_{\mathcal{V}}-\hat{y}_{\mathcal{V}}\|^2_2
+
\lambda_\mathcal{E}\|\hat{x}_{\mathcal{E}}-\hat{y}_{\mathcal{E}}\|^2_2,
\end{equation}
with $\lambda_\mathcal{V},\lambda_\mathcal{E}>0$ scaling node and edge contributions. This is the
embedding-space cost $c(x,y):=\tilde c(\hat{x},\hat{y})$, and it is differentiable in both
arguments. To compare graphs in a permutation-invariant way, we lift the local cost by optimizing
over node relabelings to obtain the \gls{fgw} cost \citep{vayer2019optimal}:
\begin{equation}\label{eq:fgw-cost}
c_{\mathrm{FGW}}(x,y)
=
\min_{\sigma\in S_n}
\left\{
c_{\mathrm{loc}}(x,\sigma\cdot y)
\right\}.
\end{equation}
This quantity is used for minibatch pairings and transport supervision. Exact evaluation is
intractable, so we use fast approximations: histogram matching for noisy--clean pairing and
node-matching for molecule--molecule pairing (Appendix~\ref{app:matching}). These pairings map each
source/noise graph $x^0$ to a data graph $x^\mathrm{data}=T(x^0)$ and define displacements
$v=\hat{x}^\mathrm{data}-\hat{x}^0$ that supervise the transport-aligned drift used to train
$-\nabla_{\hat{x}}V_\theta$ (Section~\ref{sec:training}).



\begin{figure}[h]
\centering
\resizebox{0.95\columnwidth}{!}{%
\begin{tikzpicture}[>=Latex, font=\small]

% ---------------- Panel box ----------------
\def\W{11.2}
\def\H{4.9}
\node[font=\small\bfseries] at (\W/2-0.7,\H+0.45)
{Discrete, geometry-aware energy-based proposal (local edits)};

% ---------------- Center + shells ----------------
\coordinate (xhat) at (3.8,1.55);

\pgfmathsetmacro{\rOne}{0.85}
\pgfmathsetmacro{\rTwo}{1.35}
\pgfmathsetmacro{\rThr}{1.85}

\draw[cb_lightgray, dashed, line width=0.7pt] (xhat) circle[radius=\rOne];
\draw[cb_lightgray, dashed, line width=0.7pt] (xhat) circle[radius=\rTwo];
\draw[cb_lightgray, dashed, line width=0.7pt] (xhat) circle[radius=\rThr];

\node[cb_lightgray, font=\scriptsize, align=center] at ($(xhat)+(-2.3,1.2)$)
{distance penalty\\isolines};

\node[cb_lightgray, font=\scriptsize\bfseries] at ($(xhat)+(115:\rOne)+(0.06,-0.24)$) {1 edit};
\node[cb_lightgray, font=\scriptsize\bfseries] at ($(xhat)+(115:\rTwo)+(0.06,-0.26)$) {2 edits};
\node[cb_lightgray, font=\scriptsize\bfseries] at ($(xhat)+(115:\rThr)+(0.06,-0.28)$) {3 edits};

\filldraw[black, fill=cb_red] (xhat) circle (2.2pt);
\node[font=\scriptsize, anchor=north east] at ($(xhat)+(-0.06,-0.12)$) {$\hat{x}$};
\node[cb_lightgray, font=\scriptsize\bfseries, anchor=west] at ($(xhat)+(0.12,-0.02)$) {stay};

% ---------------- Gradient direction ----------------
\pgfmathsetmacro{\gdeg}{30} % degrees
\coordinate (gend) at ($(xhat)+(\gdeg:3.05)$);
\draw[->, ultra thick, cb_green] (xhat) -- (gend);
\node[font=\small\bfseries, anchor=west, cb_green]
  at ($(xhat)+(1.8,1.0)$)
  {$-\nabla_{\hat{x}} V_\theta(\hat{x})$};

% ---------------- Dot area encodes probability ----------------
\pgfmathsetmacro{\A}{2.0}
\pgfmathsetmacro{\B}{1.0}

\pgfmathsetmacro{\sOne}{exp(\A*\rOne - \B*\rOne*\rOne)}
\pgfmathsetmacro{\sTwo}{exp(\A*\rTwo - \B*\rTwo*\rTwo)}
\pgfmathsetmacro{\sThr}{exp(\A*\rThr - \B*\rThr*\rThr)}
\pgfmathsetmacro{\smax}{max(\sOne,max(\sTwo,\sThr))}

\pgfmathsetmacro{\radMin}{0.35}
\pgfmathsetmacro{\radScale}{2.25}

\newcommand{\DrawRing}[3]{%
  \foreach \ang in {#2} {%
    \pgfmathsetmacro{\score}{exp(\A*#1*cos(\ang-\gdeg) - \B*#1*#1)}%
    \pgfmathsetmacro{\p}{\score/\smax}%
    \pgfmathsetmacro{\rad}{\radMin + \radScale*sqrt(\p)}%
    \filldraw[black, fill=#3] ($(xhat)+(\ang:#1)$) circle[radius=\rad pt];
  }%
}

\DrawRing{\rOne}{0,30,60,90,120,150,180,210,240,270,300,330}{cb_blue}
\DrawRing{\rTwo}{15,45,75,105,135,165,195,225,255,285,315,345}{cb_blue}
\DrawRing{\rThr}{0,45,90,135,180,225,270,315}{cb_blue}

% ---------------- Legend + equation (spaced higher) ----------------
\node[draw, rounded corners, fill=white, inner sep=6pt,
      font=\small, anchor=north, align=left]
  at (\W/2 -0.7,\H+0.10)
{$q_\mathrm{mixing}(x\to y)\propto \exp\bigl(
- \lambda^L_{\mathcal{V}}\|\hat{y}_{\mathcal{V}}-\hat{x}_{\mathcal{V}}\|_2^2
- \lambda^L_{\mathcal{E}}\|\hat{y}_{\mathcal{E}}-\hat{x}_{\mathcal{E}}\|_2^2$\\[4pt]
$\quad\quad\quad\quad\quad\quad - \beta^L\,\nabla V_\theta(\hat{x})^\top(\hat{y}-\hat{x})\bigr)$};

% ---------------- Dot area explanation repositioned ----------------
\node[font=\small, anchor=west] at (5.6,1.9)
{dot area $\propto\, q_\mathrm{mixing}(x\to y)$};

\end{tikzpicture}%
}
\caption{\textbf{Proposal Scoring.} Local graph edit proposals with scoring given by $q_\mathrm{mixing}(x\to y)$. The size of dots encodes proposal probabilities. Alignment with the gradient direction increases the probability, while distance penalties discourage larger edits. If the candidate is the origin (stay), we resample to promote exploration. (with equal node/edge weights $\lambda^L_{\mathcal{V}} = \lambda^L_{\mathcal{E}}$ for illustrative purposes)}
\label{fig:proposal-middle-softmax-alpha-half}
\end{figure}



\subsection{Discrete Transport-Aligned Proposal ($\epsilon\to 0$)}{\color{red} TODO: point to details}

\label{proposals}

Given an optimal coupling $\gamma$ (in practice approximated using tractable methods, e.g.,~\cite{albergo2024stochastic,tong2024improving}) and its induced transport map $T$, we assume $T$ admits a consistent representation of the form $y=T(x) \iff \hat{y}=\tilde{T}(\hat{x})$, with $x$ drawn from the source/noise distribution and $y$ from the data distribution, where $\tilde{T}$ is a continuous map that realizes $T$ under the embedding.  Expressing the JKO variational objective from \Cref{jko_em} as an optimization over transport maps (see \Cref{app:jko}) and substituting the time-step parameter $\Delta t$ with the jump step-size $\eta$ yields:
\begin{equation}
\label{eq:new_T}
    T(x)\in\arg\min_{y\in\mathcal{X}}\left\{\frac{1}{2\eta}c(x,y)+V_\theta(y)\right\},
\end{equation}
where $\pi_{k+1}=T_{\#}\pi_k$ denotes the pushforward of $\pi_k$.



Assuming the minimizer $T(x)$ lies in the interior, the first-order optimality condition becomes
\begin{equation}
0 = \left.\nabla_{\hat{y}}\left[\frac{1}{2\eta}\tilde{c}(\hat{x},\hat{y})+\tilde{V}_\theta(\hat{y})\right]\right|_{\hat{y}=\tilde{T}(\hat{x})},
\end{equation}
which implies
\begin{equation}
\nabla_{\hat{y}}\tilde V_\theta\!\left(\tilde T(\hat{x})\right)
\;=\;
-\frac{1}{2\eta}\,
\nabla_{\hat{y}}\tilde c\!\left(\hat{x},\tilde T(\hat{x})\right).
\end{equation}
This optimality condition fully characterizes the solution and can thus be directly employed as a training objective. We utilize this approach in \Cref{sec:training}, incorporating the global, permutation-invariant cost from \Cref{eq:fgw-cost}.

We are now, at sampling, trying to solve \Cref{eq:new_T} given a trained network on the above objective. We want to make the proposal local as to iterate the descent in the potential in more than one-shot such that we can also do a taylor expension on $V(y)$ around x. Rather than controlling the step-size parameter $\eta$, we simplify the (local) greedy proposal by restricting the kernel to exactly $N$ edits per jump. For this proposal we consider the cost to be $c_{\mathrm{loc}}(x,y)$ with $\lambda_\mathcal{V}=\lambda_\mathcal{E}$ which yields fixed costs for every such candidate - which means the cost is constant and can be neglected within the search. Linearizing $V_\theta(y)\approx V_\theta(x)+\Delta\hat{x}^\top\nabla_{\hat{x}} V_\theta(\hat{x})$ around $\hat{x}$, with $\Delta\hat{x}=\hat{y}-\hat{x}$, we deterministically select the top-$N$ candidate edits ranked by steepest descent. Formally, the candidate set $C_N(x)$ is defined as the set of all edits achievable by exactly $N$ discrete modifications (changing/removing edges or changing nodes) starting from state $x$. The greedy  proposal is thus given by:
\begin{gather}
y^*(x) = x + \arg\min_{\Delta x \in C_N(x)} (\Delta\hat{x})^\top\nabla_{\hat{x}} V_\theta(\hat{x}), \\
q_{\text{greedy}}(x\to y) = \delta_{y^*(x)}(y).
\end{gather}

In the transport regime, we project the proposed moves onto directions with strictly negative gradient, rejecting any edits that do not satisfy this condition. It is deterministic. If no candidate edits yield a negative gradient direction, the greedy sampler becomes trapped, prompting a transition to the Langevin-based proposal $q_\mathrm{mixing}$.

\subsection{Discrete Mixing Proposal ($\epsilon>0$).}
Discretizing the dynamics (\Cref{em_sde}) in the embedding space  with temperature $\epsilon$, step size $\eta$, and graphs $x,y \in \mathcal{X}$ yields:
\begin{equation}
q_{\mathrm{mixing}}(x\to y)
\,\propto\,
\exp\left(
-\frac{1}{4\epsilon\eta}\|\hat{y}-\hat{x}+\eta\nabla V_\theta(\hat{x})\|^2
\right).
\end{equation}

Expanding explicitly and dropping terms independent of $\hat{y}$:
\begin{align}
\log q_\mathrm{mixing}(x\to y)
&=
-\frac{1}{4\epsilon\eta}\|\hat{y}-\hat{x}\|^2
- \frac{1}{2\epsilon}\nabla V_\theta(\hat{x})^\top(\hat{y}-\hat{x})\notag\\
&\quad+\text{const.}
\end{align}
The proposal is effectively local because of the distance penalty, although technically it can still propose an arbitrary number of edits. We define Langevin-specific parameters $(\lambda_{\mathcal{V}}^L, \lambda_{\mathcal{E}}^L, \beta^L)$ absorbing all constants (including $\epsilon$, $\eta$, and numerical factors):
\addtocounter{equation}{1}
\begin{align}
&q_\mathrm{mixing}(x\to y)
\;\propto\tag{\theequation}\\
&\quad\exp\!\bigl(
- \lambda_{\mathcal{V}}^L\|\hat{y}_{\mathcal{V}}-\hat{x}_{\mathcal{V}}\|_2^2
- \lambda_{\mathcal{E}}^L\|\hat{y}_{\mathcal{E}}-\hat{x}_{\mathcal{E}}\|_2^2
- \beta^L\,\nabla V_\theta(\hat{x})^\top(\hat{y}-\hat{x})
\bigr).\notag
\end{align}


The resulting kernel can be viewed as a \emph{linearized locally balanced proposal}; related
discrete-gradient samplers appear in~\citep{zanella2020informed, grathwohl2021oops,zhang2022langevin} for particular costs or candidate sets.
Figure~\ref{fig:proposal-middle-softmax-alpha-half} illustrates the resulting geometry-aware proposal
scoring and how gradient alignment competes with distance penalties.

In practice, this Langevin-based proposal is combined with a Metropolis--Hastings acceptance step, ensuring convergence to the target distribution $\pi_\theta(x)\propto\exp(-\beta_{mh} V_\theta(x))$. 

We calibrate the parameters $\beta_{mh}$, $\beta^{L}$, $\lambda_{\mathcal{V}}^{L}$, and $\lambda_{\mathcal{E}}^{L}$ by maximizing the relative likelihood of generated samples compared to the training dataset samples under the learned energy $V_\theta$, directly after the warmup training phase described in \Cref{sec:training}.



In the mixing regime, we accept the stochastic proposal with \gls{mh} probability:
\begin{equation}\label{eq:mh-accept}
\begin{aligned}
\alpha(x,y)
&=\min\left(1,\frac{\pi_\theta(y)\,q_{\mathrm{mixing}}(y\to x)}{\pi_\theta(x)\,q_{\mathrm{mixing}}(x\to y)}\right).\\
\end{aligned}
\end{equation}
This choice
enforces detailed balance:
\begin{equation}\label{eq:mh-detailed-balance}
\begin{aligned}
\pi_\theta(x)\,q_{\mathrm{mixing}}(x\to y)\,\alpha(x,y)
= \\
\pi_\theta(y)\,q_{\mathrm{mixing}}(y\to x)\,\alpha(y,x)
\end{aligned}
\end{equation}
so $\pi_\theta$ is stationary for the mixing dynamics~\citep{metropolis1953equation,hastings1970monte}.
Implementation details of the proposal kernel (factorized logits, temperature annealing, and do-nothing
resampling) are provided in Appendix~\ref{app:sampling-details}.



\subsection{Initializations and proposal schedules.}

\paragraph{Determining the Regime}
\label{switch}
We control the regime through an energy-based binary indicator $s\in\{0,1\}$: $s=0$ selects transport (deterministic $\epsilon\to 0$ limit), while $s=1$ selects mixing (finite $\epsilon>0$). If initialized at noise, the sampler follows deterministic greedy edits guided strictly by negative gradient directions. Transition from transport to mixing occurs when the Markov chain either reaches the target energy (defined as the mean data energy estimated during training) or becomes trapped in a local minimum—no valid gradient-improving edits exist—for all samples, whichever happens first. Once triggered, the sampler enters the stochastic Langevin-based mixing regime. More about proposal schedules in \Cref{init_prop}.

\paragraph{Init.... TODO}
\label{init_prop}
We use two initialization regimes with aligned proposal schedules:
\begin{itemize}[leftmargin=*]
\item \textbf{Noise initialization:} sample the node count from the empirical histogram, then draw node/edge
types uniformly (uniform noise) or from empirical marginals (marginal noise). These start far from
the data distribution, so we run transport with greedy proposals $q_{\text{greedy}}$ until the energy-based switch (see
\Cref{switch}), then enter mixing with the proposal
$q_{\mathrm{mixing}}$. 
\item \textbf{Data initialization:} draw graphs from the training set. We skip transport and start directly in mixing with $q_{\mathrm{mixing}}$, using a higher initial proposal temperature that is annealed before
fixed-temperature mixing (Appendix~\ref{app:sampling-details}).
\end{itemize}
All initializations start near-zero \gls{vun}. Figure~\ref{fig:moses-energy-and-metrics}
summarizes the resulting energy trajectories under noise vs.\ data initialization and the
corresponding proposal regimes.

\begin{figure}[h]
\centering

\begin{tikzpicture}[scale=1.15]
  \draw[->, thick] (0,-0.36) -- (6.2,-0.36);
  \draw[->, thick] (0,-0.36) -- (0,2.6);
  \node[rotate=90] at (-0.7,1.12) {$V_\theta(x)$};
  \node[font=\small\bfseries] at (3.1,2.9) {Energy Trajectories};

  % Axes ticks
  \foreach \x/\lab in {0/0,1.2/100,2.4/200,3.6/300,4.8/400,6.0/500} {
    \draw (\x,-0.36) -- (\x,-0.42);
    \node[below, font=\scriptsize] at (\x,-0.42) {\lab};
  }
  \foreach \y/\lab in {-0.36/-40,0/0,0.9/100,1.8/200,2.475/275} {
    \draw (0,\y) -- (-0.08,\y);
    \node[left, font=\scriptsize] at (-0.1,\y) {\lab};
  }

  % Annealing window (S) - Using TEAL (was Green)
  \draw[decorate, decoration={brace, amplitude=4pt, mirror}, draw=cb_teal] (0.06,0.2) -- (2.10,0.2);
  \node[font=\scriptsize, text=cb_teal] at (1.36,-0.16) {annealed temperature mixing};
  
  % Fixed-temperature window - Using BLUE (was Blue)
  \draw[decorate, decoration={brace, amplitude=4pt, mirror}, draw=cb_blue] (2.28,0.2) -- (5.94,0.2);
  \node[font=\scriptsize, text=cb_blue] at (4.20,-0.16) {fixed-temperature mixing};

  % Bands (range)
    % Noise band (Orange)
    \fill[cb_orange!20] (0,1.78) rectangle (6.0,2.28);
    \draw[cb_orange, line width=0.6pt] (0,2.03) -- (6.0,2.03);
    
    % Data band (Blue)
    \fill[cb_blue!20, opacity=0.5] (0,0.18) rectangle (6.0,0.68);
    \draw[cb_blue, line width=0.6pt] (0,0.43) -- (6.0,0.43);


  % Greedy proposal window - Using RED (was Red)
  \draw[decorate, decoration={brace, amplitude=4pt}, draw=cb_red] (0.06,2.34) -- (2.22,2.34);
  \node[font=\scriptsize, text=cb_red] at (1.14,2.65) {greedy proposals};

  % Trajectories (means + std)
  
  % Noise Initialized (RED)
  \fill[cb_red!20, opacity=0.6] plot[smooth] coordinates {
    (0.0,2.392) (0.7,1.992) (1.3,1.492) (2.1,0.742)
    (2.7,0.732) (3.3,0.732) (4.2,0.732) (5.1,0.732) (6.0,0.732)
  } -- plot[smooth] coordinates {
    (6.0,0.172) (5.1,0.172) (4.2,0.172) (3.3,0.172)
    (2.7,0.172) (2.1,0.182) (1.3,0.932) (0.7,1.432) (0.0,1.832)
  } -- cycle;
  \draw[cb_red, densely dotted, line width=1.1pt] plot[smooth] coordinates {
    (0.0,2.112) (0.7,1.712) (1.3,1.212) (2.1,0.462)
    (2.7,0.452) (3.3,0.452) (4.2,0.452) (5.1,0.452) (6.0,0.452)
  };

  % Data Initialized (TEAL - was Green)
  \fill[cb_teal!30, opacity=0.5] plot[smooth] coordinates {
    (0.0,0.74) (0.4,0.91) (0.9,0.82) (1.4,0.74)
    (2.2,0.72) (3.1,0.71) (4.2,0.71) (5.1,0.71) (6.0,0.71)
  } -- plot[smooth] coordinates {
    (6.0,0.18) (5.1,0.18) (4.2,0.18) (3.1,0.18)
    (2.2,0.18) (1.4,0.18) (0.9,0.26) (0.4,0.35) (0.0,0.18)
  } -- cycle;
  \draw[cb_teal, densely dotted, line width=1.1pt] plot[smooth] coordinates {
    (0.0,0.46) (0.4,0.63) (0.9,0.54) (1.4,0.46)
    (2.2,0.44) (3.1,0.43) (4.2,0.43) (5.1,0.43) (6.0,0.43)
  };

  % Legend
\begin{scope}[shift={(3.3,1.7)}]
  % Noise band (orange): band + constant (solid) reference line
  \draw[fill=cb_orange!20, draw=cb_orange] (0,0) rectangle (0.35,0.18);
  \draw[cb_orange, line width=0.6pt] (0.04,0.09) -- (0.31,0.09);
  \node[anchor=west, font=\scriptsize] at (0.45,0.09) {noise band};

  % Data band (blue): band + constant (solid) reference line
  \draw[fill=cb_blue!20, draw=cb_blue, fill opacity=0.5] (0,-0.28) rectangle (0.35,-0.10);
  \draw[cb_blue, line width=0.6pt] (0.04,-0.19) -- (0.31,-0.19);
  \node[anchor=west, font=\scriptsize] at (0.45,-0.19) {data band};

  % Noise initialized samples (red): band + dotted mean line
  \draw[fill=cb_red!20, draw=cb_red] (0,-0.56) rectangle (0.35,-0.38);
  \draw[cb_red, densely dotted, line width=1.1pt] (0.04,-0.47) -- (0.31,-0.47);
  \node[anchor=west, font=\scriptsize] at (0.45,-0.47) {samples: noise initialized};

  % Data initialized samples (teal): band + dotted mean line
  \draw[fill=cb_teal!30, draw=cb_teal] (0,-0.84) rectangle (0.35,-0.66);
  \draw[cb_teal, densely dotted, line width=1.1pt] (0.04,-0.75) -- (0.31,-0.75);
  \node[anchor=west, font=\scriptsize] at (0.45,-0.75) {samples: data initialized};
\end{scope}

\end{tikzpicture}

\vspace{0.25cm}

\begin{tikzpicture}[xscale=0.99, yscale=1.06]
  \draw[->, thick] (0,0) -- (7.2,0) node[midway, below=14pt] {Inference \glsentryshort{mcmc} steps};
  \draw[->, thick] (0,0) -- (0,3.35);
  \node[rotate=90] at (-0.9,1.6) {Fraction$\uparrow$};

  % Title
  \node[font=\small\bfseries] at (3.5,3.55) {Validity \& Novelty Trajectories};

  % Axes ticks
  \foreach \x/\lab in {0/0,1.4/100,2.8/200,4.2/300,5.6/400,7.0/500} {
    \draw (\x,0) -- (\x,-0.06);
    \node[below, font=\scriptsize] at (\x,-0.06) {\lab};
  }
  \foreach \y/\lab in {0/0.00,0.8/0.25,1.6/0.50,2.4/0.75,2.88/0.90,3.2/1.00} {
    \draw (0,\y) -- (-0.08,\y);
    \node[left, font=\scriptsize] at (-0.1,\y) {\lab};
  }

  % Validity (Data Init) - Using TEAL
  \draw[cb_teal, thick] plot[smooth] coordinates {
    (0.0,3.20) (0.1,1.40) (0.2,0.55) (0.3,0.26) (0.6,0.22)
    (1.0,0.30) (1.4,0.42) (1.8,0.70) (2.2,1.05) (2.6,1.40)
    (3.0,1.70) (3.4,2.00) (3.8,2.25) (4.2,2.45) (4.6,2.60)
    (5.0,2.72) (5.6,2.86) (6.2,3.06) (6.8,3.075) (7.0,3.088)
  };

  % Novelty (Data Init) - Using TEAL (Dashed)
  \draw[cb_teal, thick, dashed] plot[smooth] coordinates {
    (0.0,0.00) (0.1,0.05) (0.2,0.25) (0.3,0.80) (0.4,1.60)
    (0.5,2.72) (1.0,2.45) (1.5,1.70) (2.0,2.05) (2.5,2.30)
    (3.0,2.45) (3.5,2.60) (4.0,2.70) (4.5,2.78) (5.0,2.85)
    (5.5,2.90) (6.0,2.95) (6.5,2.98) (7.0,3.00)
  };

  % Validity (Noise Init) - Using RED
  \draw[cb_red, thick] plot[smooth] coordinates {
    (0.0,0.00) (0.6,0.08) (1.0,0.32) (1.4,1.12) (1.8,2.24)
    (2.2,2.60) (2.6,2.72) (3.0,2.80) (4.2,2.88) (5.6,2.92) (7.0,2.944)
  };
  % Novelty (Noise Init) - Using RED (Dashed)
  \draw[cb_red, thick, dashed] plot[smooth] coordinates {
    (0.0,3.20) (1.4,3.187) (2.8,3.174) (4.2,3.162) (5.6,3.149) (7.0,3.136)
  };

  % Legend
  \begin{scope}[shift={(3.3,0.15)}]
    \draw[fill=white, draw=cb_gray] (0,0) rectangle (3.6,1.4);
    
    % Noise Init (Red)
    \draw[cb_red, thick] (0.15,1.15) -- (0.6,1.15);
    \node[anchor=west, font=\scriptsize] at (0.7,1.15) {Validity (noise initialized)};
    \draw[cb_red, thick, dashed] (0.15,0.85) -- (0.6,0.85);
    \node[anchor=west, font=\scriptsize] at (0.7,0.85) {Novelty (noise initialized)};
    
    % Data Init (Teal)
    \draw[cb_teal, thick] (0.15,0.45) -- (0.6,0.45);
    \node[anchor=west, font=\scriptsize] at (0.7,0.45) {Validity (data initialized)};
    \draw[cb_teal, thick, dashed] (0.15,0.15) -- (0.6,0.15);
    \node[anchor=west, font=\scriptsize] at (0.7,0.15) {Novelty (data initialized)};
  \end{scope}
\end{tikzpicture}
\caption{\footnotesize \textbf{Energy and Sampling Trajectories.} Top: energy evolution for noise- vs.\ data-initialized chains, with mean $\pm$ 1 std bands and reference levels at $\mathbb{E}_{x\sim \pi_\mathrm{data}}[V_\theta(x)]$ and $\mathbb{E}_{x\sim \pi_0}[V_\theta(x)]$. Noise-initialized chains use greedy proposals to reach the data distribution; data-initialized chains use temperature annealing (low initial $\beta_{\mathrm{mh}}$, gradually increasing) to recover novelty.
 Bottom: validity and novelty trajectories for data-initialized (teal) and noise-initialized (red) chains over inference steps. Uniqueness $\approx$100\%.}
\label{fig:moses-energy-and-metrics}
\vspace{-5pt}
\end{figure}
